\chapter{Методика}
\label{ch:methods}

Диффузией  называют самопроизвольное взаимное проникновение веществ друг в друга, происходящее вследствие хаотичного теплового движения молекул. При перемешивании молекул разного сорта говорят о взаимной (или концентрационной) диффузии.

Диффузия в системе, состоящей из двух компонентов $a$ и $b$ (бинарная смесь), подчиняется закону Фика: плотности потока компонентов $j_{a ,b}$ (количество частиц, пересекающих единичную площадку в единицу времени) пропорциональны градиентам их концентраций $\nabla n_{a , b}$, что в одномерном случае можно записать как

\[j_a = -D \frac{\partial n_a}{\partial x}, \hspace{20pt} j_b = -D \frac{\partial n_b}{\partial x},\]

где $D$ — коэффициент взаимной диффузии компонентов. Знак «минус» отражает тот факт, что диффузия идёт в направлении выравнивания концентраций. Равновесие достигается при равномерном распределении вещества по объёму сосуда ($ \partial n /\partial x = 0$ ).
